
\chapter{Subject focus in Spanish and French: an unsolved matter}\label{sec:2}
\section{Introduction}

The argument asymmetry described in Chapter 1 is particularly visible in the realization of focus in Spanish and French. Both languages, in fact, are said to adopt non-canonical means to focus the subject, while a canonical SVO order is generally preferred to express object focus (cf. \cite{lambrecht1996information} for both languages). However, there are considerable differences between the two languages. In many theoretical views, focused subjects in Spanish are associated with the re-ordering of the constituents in the sentence, where the focused subject appears in a postverbal position, rather than its usual preverbal (topical) one. 

\ea \label{ex526}
\ea \textit{¿Quién abrió la puerta?}\\
 ‘Who opened the door?'
\ex \textit{La abrió} {[\textit{Maria}]\textsubscript{\textsc{foc}}}\normalsize.\\
`Maria opened it.'
\z \z 

This word order, generally called postverbal, VS order or subject-verb inversion, is also reported for the majority of members of the Romance family. Although micro-differences can be observed, Italian, Romanian, Catalan, European Portuguese and many non-standard varieties make use of this possibility to mark the subject as narrow focus of the sentence.

In the case of French, however, postverbal subjects (see \cite{lahousse2006} and \textcite{Lahousse2011} for an exhaustive overview of the VS order in French) do not exhibit the same information-structural properties of their Spanish counterpart, and are almost exclusively found in written (mainly literary) texts \parencite{delbecque2006estudio}. \textcite{dufter201614}, for instance, claim that the possibility of VS in French is restricted to very few cases:  it is felicitous with unaccusative verbs and in passive constructions. However, it is marginally possible with locative inversion and, notably, severely constrained in transitive constructions, such as in (\ref{ex25}): 

 \ea \label{ex25}
  \textit{Rendront un devoir} {[\textit{les élèves qui ont raté l’examen de chimie}]\textsubscript{\textsc{foc}}}\normalsize.\\
 `The students that failed the chemistry exam must submit an assignment.' \hfill \parencite[431]{dufter201614}    
\z

According to \textcite{lambrecht1996information}, in spoken French, the corresponding means to focus the subject is a \textit{c'est} cleft. 

 \ea \label{ex527}
\ea \textit{Qui a ouvert la porte?}\\
 ‘Who opened the door?'
\ex \textit{C'est}  {[\textit{Marie}]\textsubscript{\textsc{foc}}} \normalsize \textit{qui l'a ouverte}.\\
`It is Marie who opened it.'
\z \z 

This difference between the French and the Spanish (or most of the other Romance languages) focus-marking system is generally attributed to the difference in the syntactic and prosodic characteristics of the two languages. As is generally acknowledged, the possibility of realizing the subject in postverbal position (also called ``free inversion”) is part of the ``Null Subject Parameter” (cf. \cite{rizzi1982issues, biberauer2009parametric}, \textit{inter alia}), and the VS order has often been considered a major strategy for subject focusing in a pro-drop language  like  Italian or Spanish (cf. \cite{belletti2008answering}). It should be noted, however, that the existence of correlations between the positive setting of the Null Subject Parameter and other grammatical properties—such as the possibility of free inversion—has been repeatedly questioned \parencite{Newmeyer2005}. In particular, typological research \parencite{Gilligan1987} has shown that the purported correlation between null subjecthood and freedom of constituent order is empirically untenable.

Nevertheless, the lack of postverbal subjects and the prevalent use of clefts in French is, at least in part, attributed to the fact that this language has both a relatively rigid constituent order and a relatively rigid rhythmic structure \parencite{dufter2003rhythmic}.
According to \textcite{dufter201614}, deviations from the basic declarative subject-verb-object order have become less and less available over time in the history of French.

Moreover, according to \textcite{lambrecht2001framework}, ``the occurrence of cleft constructions in a language correlates with the degree of positional freedom of prosodic accents and syntactic constituents in that language” \parencite[488]{lambrecht2001framework}. 
Prosodically, Spanish is a language that displays lexical stress. This implies that speakers can signal focus by assigning pitch accent to a particular constituent. In French, on the other hand, prosodic  prominence   is  obligatorily realized at the right edge of the Phonological Phrase, and cannot be displaced from this default position in order to signal focus on non-phrase-final constituents \textit{in situ}. Thus, French does not display pitch accents to mark focus, but rather assigns prosodic prominence at the right edge of the phonological phrase through phrasing, i.e., a boundary tone, and not a pitch accent \parencite{fery2001focus}.

In sum, while Spanish can assign prosodic prominence on a particular constituent, French needs to create phrasing. A cleft sentence would therefore be an optimal strategy to mark focus prominence, by creating two intonational phrases (through a copular sentence and a relative clause) and by positioning the subject at the end of that phonological phrase, i.e., at the boundary tone, so that it can receive prosodic prominence. This would explain why, in the case of object focus, no movement of the constituents or specific construction is needed. SVO being the canonical word order in both French and Spanish, a direct object is already in a position where it can receive prosodic prominence. A canonical SVO word order is also more economic than a cleft sentence, whose bi-clausal structure makes it ``heavy” and ``costly” (cf. \cite{destruel2016focus} for an experimental study).

Thus, the usage of clefts is generally considered a “compensatory device” (\cite[457]{doherty2001discourse}; \cite[487]{di2006clefting}) for the lack of word order flexibility and lexical stress, and is avoided  whenever alternative focus-marking strategies are possible. Such a ``trade-off” view on clefts, however, has been questioned by recent diachronic studies \parencite{dufter2008explaining} that have investigated the frequency of clefts from Old French, a null-subject language with lexical stress different from Modern French, which, as mentioned above, has become very restricted in terms of word order flexibility and accentuation. The results of this diachronic study have revealed an overuse of clefts even in the presence of other, more economical focus-marking devices, suggesting that the rise and the frequent usage of cleft constructions in Modern French must be attributed to factors other than this ``compensatory” function.

\textcite{lambrecht2001framework} considers both the Spanish VS order and the French \textit{c'est} clefts as operations of ``detopicalization” of the subject, thus advancing the claim that these two strategies, even if syntactically considerably different, share important pragmatic similarities. This assumption has been repeatedly challenged by several empirical studies, that have shown that variation in the realization of focus can be observed not only cross-linguistically but also language-internally. Indeed, in both Spanish and French, experimental studies have revealed the possibility of \textit{in situ} focus, even in non-contrastive interpretations (e.g., \cite{gabriel2007fokus, gabriel2010focus, heidinger2013information, bosch2013variation}, a.o.). Moreover, recent corpus studies have found out that different cleft types (such as \textit{il y a} clefts) can have the same IS-partition of \textit{c'est} clefts, thus competing for the same pragmatic function \parencite{karssenberg2017oiseaux, karssenberg20188}.

In this chapter, the many contributions to the investigation of subject focus are reviewed and discussed extensively. I present studies conducted by both formalist and functionalist authors, highlighting the main findings and problematic issues. The chapter is organized as follows: in section 2.1, I provide an overview of the studies on subject focus in Spanish, both from a theoretical and an empirical perspective. Section 2.2 summarizes the theoretical and empirical studies conducted for subject focus in French. In section 2.3, I compare the factors that have been claimed to interact with focus realization in the two languages. Finally, in section 2.4, I point out the research gap that I intend to fill with the present investigation.


\section{Subject focus in Spanish: previous studies}


Concerning the question of the realization of Spanish focused subjects, it is possible to distinguish two schools of thought:  on one side, the theoretical approaches (cf. \cite{zubizarreta1998prosody, ramalle2005manual, costa2001postverbal, gutierrez2013structural, samek2001crosslinguistic}, \textit{inter alia}), on the other side, the functional ones (cf. \cite{gabriel2007fokus, gabriel2010focus, heidinger2013information, uth2014spanish, hualde2002intonation}, \textit{inter alia}).

In a nutshell, in theoretical approaches, focused subjects in Spanish must necessarily occur in IP-final position, making postverbal subjects (or VS order) the only possible way to realize information focus on the subject. 

 \ea \label{ex501} \textit{Comió una manzana} {[\textit{Juan}]\textsubscript{\textsc{foc}}}\normalsize. \\
`Juan ate an apple.'
\z


In this view, \textit{in situ} (or prosodic) marking of focus on the subject is reserved to a contrastive interpretation of the subject referent \parencite[4228]{zubizarreta1999word}. 
However, such claims are often based on introspection or grammaticality judgements. The functional side does not entirely contradict the position defended by formal approaches. Rather, it advocates a weaker version of the above theory, showing through experiments that subject focus can also be realized by other strategies, such as \textit{in situ} focusing \parencite{gabriel2007fokus, heidinger2013information, uth2014spanish} or clefts \parencite{feldhausen2015oraciones}, even in the case of a non-contrastive reading.

The main distinction between the two views can be resumed as such: while formal approaches argue in favour of a deterministic relationship between form and function, functional studies acknowledge the existence of variation and the possibility for the same pragmatic function to be expressed by different forms. I will first review the formal schools and then the functional ones, pointing out the factors that have been claimed to determine the optionality among different focus-marking devices, and which are relevant for my argumentation.

\subsection{The formal framework: postverbal subjects}
 
The term postverbal subject or VS order covers a series of syntactic phenomena that are not all necessarily related to focus. Subject-verb inversion can occur in a range of different linguistic contexts that are either constrained by syntactic requirements \footnote{For instance, obligatory subject verb inversion with focus fronting (FF) and in wh-questions (see \cite{gutierrez2013structural}).} or justified by stylistics reasons. Since this book is concerned with focus, the above-mentioned cases are not discussed here. Additionally, cases where the VS order expresses an all-focus interpretation or thetic judgement \parencite{sasse1987thetic} are also out of the scope of the present study. Rather, postverbal subjects are considered only when they function as a focalizing means for the subject alone. There are various explanations that can be found in the literature of why a focused subject must appear sentence-finally in Spanish:


\begin{description}
    \item[1. \textbf{Prosodic reasons}]
    
\textcite{cinque1993null} relates focus to stress and stress assignment rules.  According to his view, stress must be assigned to the most embedded metrical constituent, which also corresponds to a syntactic constituent. He proposes that the accent indicating focus in a phrase must coincide with the most prominent accent (nuclear accent) in the phrase. This implies that the focus accent must fall on the most embedded constituent of the phrase as well. In languages that exhibit canonical SVO order, like Spanish, the most embedded constituent is the object.\footnote {The term \textit{object} here is an umbrella term that includes direct objects, indirect objects or prepositional phrase, depending on the argument structure of the verb.} Thus, in the case of object focus, no conflict can be observed between nuclear stress and focus accent, since it is always the rightmost element that receives the nuclear stress. Conversely, in the case of subject focus, intonation, constituent order requirements and focus prominence can be in conflict with one another \parencite{selkirk1995sentence, truckenbrodt1999relation}.

\textcite{truckenbrodt1999relation}, for instance, has shown that sacrificing the canonical pattern of phonological phrasing is a strategy that some languages use to solve such potential conflicts.
\textcite{zubizarreta1998prosody}, in her seminal work on postverbal subjects, suggested that in order to assign stress to the most deeply embedded element even under information focus, languages use two options: (i) render non-focused material invisible for the stress assignment algorithm, so that the focused constituent automatically qualifies as the most deeply embedded constituent; or (ii) alter word order to get the focused XP to be the most deeply embedded item. English has been argued to belong to the first category, whereas Spanish has been attributed to the latter class. Thus, in Zubizarreta's Minimalist approach, the  object  would  undergo p(rosodically-motivated) movement to allow the focused subject to be the most deeply embedded element and, thus, to receive prosodic prominence (but see \cite{ordonez2000clausal}, a.o., for an alternative view, positing the existence of object scrambling past the \textit{in situ} subject). Movement operations such as these would be triggered by abstract IS-features situated in the underlying structure of the sentence (e.g., [Focus]).

P(rosodically motivated) movement is assumed to be obligatory in the case of information focus, and any deviation from this pattern is considered pragmatically infelicitous or even ungrammatical. This is also the view shared by classical optimality-theory (OT) approaches (e.g.,  \cite{costa2001postverbal, gutierrez2013structural, samek2001crosslinguistic}), according to which constructions with non-final focus are ruled out due to a violation of a high-ranked constraint such as the NSR (“Nuclear stress rule”), requiring nuclear stress to fall on the rightmost constituent of the sentence. A SVO configuration, with the subject receiving prosodic prominence would thus be acceptable only in the case of a contrastive interpretation.

Building on these observations, \textcite{buring2001focus} suggest an analysis where the focus and intonational requirements of English, Spanish and German are essentially the same, contrary to what has been suggested in previous work.  Instead, these languages differ in when and whether they sacrifice phonological phrasing requirements or syntactic requirements.  As such, their analysis preserves the insights captured in \textcite{zubizarreta1998prosody}, while maintaining the assumption, shared in most recent work on phonological phrasing and accenting, that prominence is not governed by parameterized syntax-based principles such as the Nuclear Stress Rule. Instead, they propose that movement of focused XPs is the strategy used to accommodate these conflicts in languages where sacrificing the canonical pattern of phonological phrasing is not an option. The difference between focus marking between, say, English and Spanish, is not prosodic prominence versus syntactic position, but rather prosodic prominence \textit{in situ} versus prosodic prominence in a specific position.\\

\item[2.  \textbf{Syntactic reasons}]

Syntactic accounts of focus (starting from \cite{brody1990some}) encode focus  syntactically  by  proposing  the  existence  of  a  ‘Focus  Criterion’, according to which a focused item must move to the spec of a FocusPhrase (FocP) to enter into a Spec/Head relation with a head endowed with a Focus feature. A significant work of this type is proposed by \textcite{rizzi1997fine}, in his Cartographic Approach, that divides the left periphery into a series of functional heads which are represented below.

\ea{}
[ForceP [ TopP*[ \textbf{FocP} [ TopP* [ FinP [ IP ...
\z

According to Rizzi’s theory, preverbal focus is situated in a designated functional projection (FocP) within the split-CP (or left periphery of the sentence) above the IP. While his theory accounts for the behavior of preverbal focus, it fails to provide an explanation for the postverbal one. \textcite{belletti2001inversion, belletti2004aspects, belletti2008cp, belletti2008answering} integrate \textcite{rizzi1997fine} into a new account for postverbal focus, suggesting that  the  area  immediately  above  the  VP  (``the  lowIP area”, in her terms) is parallel to the left periphery in hosting focus and topic heads. According to the author, 
the VS order would thus be derived via movement of the subject to a dedicated position in the so-called low periphery of the sentence (cf.  also \cite{cardinaletti2004toward}).
Specifically, the postverbal subject moves to this low peripheral focus position where it is interpreted, while a silent \textit{pro} fills the preverbal subject position (meeting EPP requirements). In other words, Belletti proposes the existence of another TopP – FocP - TopP sequence just above the VP. This means that, in the spirit of her analysis, postposed subjects in such cases are moved subjects, and not subjects in their \textit{in situ} origin (as e.g., \cite{burzio1986italian} claimed).\\
Such strictly syntactic views, however, have often been the object of criticism. Many authors have pointed out that there are empirical problems with the predictions made by these analyses \parencite{samek2009topic}. In particular, it has been claimed that these syntactic accounts do not explain why focus allows for so many distinct positions. The hypothesis that contrastive focus is restricted to a CP-level projection and information focus to a TP-level \parencite{belletti2004aspects} may at first appear to address this issue, but it is inconsistent with empirical observations. Under this  hypothesis, the  distribution  of  each  focus type  should  be restricted to a single position. However, as will be discussed in the next section, both contrastive and information focus can occur in multiple positions. Additionally, the same theory predicts that the distributions of contrastive and new information focus do not overlap, but empirical studies have shown that they actually can (e.g., \cite{brunetti2009italian}).\\

\item[3. \textbf{Pragmatic reasons}]

The pragmatic explanation does not contradict the other views. Rather, it adds a communicative ``flavor” to the questions of postverbal subjects by assuming that the default reading of a subject, when it appears in a canonical preverbal position in SVO languages, is the topic one. This view was mainly proposed by \textcite{lambrecht1996information, lambrecht2010constraints}. Since preverbal subjects are generally interpreted as topics, languages possess a series of syntactic, prosodic and morphological means to avoid such interpretation. Starting from the assumption that there is a cross-linguistic correlation between certain IS-functions and grammatical roles, where subjects are usually topics, while the VP is normally interpreted as focus \parencite{drubig2003toward}, it has been proposed that speakers would operate a process of “demotion”,  “objectification” or “detopicalization” of the subject, by situating it in a postverbal position (a position generally reserved to objects). Thus, the VS order in Spanish would be an iconic means to signal to the hearer that constituent must not be interpreted as topic, but rather as the relevant part of the answer, i.e., the focus.\\
This view is also supported by \textcite{fiedler2006subject}, who, building on Romance  data,  as  well  as  Chadic  data,  suggest  that  there  is  a  strong requirement for all focus constituents, and not only non-subjects, to be placed in a prototypical, and in some sense prominent, focus position behind the verb.  Subject-inverting in the languages that allow it would thus obey the communicative requirement to  mark  the  subject  not only as  a non-topic, but especially as focus of the sentence.

\end{description}


 
 \subsection{The functional approaches: alternative focus-marking strategies}
 
 The claim that focused subjects in Spanish can be expressed only in IP-final position, while preverbal focused subjects can only be allowed with a contrastive interpretation (see section 2.2.1) has been called into question by more recent empirical studies. These studies, using methods such as reading tasks, acceptability judgements, elicitation tasks to corpus studies, have revealed that there is no one-to-one mapping between form and function, and that different syntactic positions can correspond to the same pragmatic function or, conversely, that the same pragmatic function can be expressed by different syntactic configurations (e.g., \cite{hualde2002intonation, feldhausen2014prosody, heidinger2013information, , gabriel2007fokus} and \cite{gabriel2010focus} for Argentinean Spanish; \cite{muntendam2013nature} for Andean Spanish; \cite{hoot2016narrow} for Mexican Spanish; \cite{bosch2013variation} for Canarian Spanish). 
 
 In particular, it has been shown that preverbal subjects can receive prosodic prominence even in the case of non-contrastive focus, and that contrastive focus can also be expressed by a subject appearing postverbally. \textcite{brunetti2009italian}, for instance, presents a qualitative corpus study on focus fronting in Italian and Spanish. The analysis of the contexts in which these utterances occur shows that only a subset of the fronted foci relates to an antecedent in discourse that is contrasted or corrected by the focused referent.
 
As such, empirical approaches to subject focus in Spanish do not contradict entirely the claims advanced by formal theories. Rather, they contribute to a contextualization of such claims, by showing that language-internal variation, and hence optionality, must be accounted for, instead of simply being labeled as “differences at the spell-out”, “competing grammars” or “dialectal differences”. \textcite{alonso2002null}, for instance, reports two experiments in which 12 question–answer pairs were submitted to 80 and 72 participants, respectively. Two different potential answers were offered for each question,  one with SV order (\ref{ex26}b) and one with VS order (\ref{ex26}b').


 \ea \label{ex26}
 \ea \textit{Quién vino?}\\
 `Who came?'
\ex {[\textit{Juan}]\textsubscript{\textsc{foc}}} \normalsize \textit{vino}.\\
`Juan came.'\\
or
\exi{b\small{$'$}.} \textit{Vino} {[\textit{Juan}]\textsubscript{\textsc{foc}}}\normalsize.\\
`Juan came.' \hfill \parencite[17]{alonso2002null}    
\z \z


Contrary to the predictions of a theory like the one proposed by \textcite{zubizarreta1999word}, the answer \textit{Juan vino} (\ref{ex26}b), in which the subject is the first constituent, and thus does not receive default stress, is chosen as the answer to the question `Who came?' 51.5\% of the time (as reported by \cite{biezma2014multiple}).  Important additional evidence is provided by \textcite{ocampo2003notion}. Building on spontaneous conversations, the author shows that main intonational prominence does not necessarily fall on the non-presupposed part of the sentence, i.e., the focused part of the sentence, and that speakers of Spanish used stress for many different functions.  In the author's analysis, the part of the sentence presenting new information does not actually bear stress.
The results from his corpora add to the evidence that the subject in non-marked word order (SVO) can be conveyed as focused, regardless of default intonational prominence on the object. 

Another prominent empirical contribution, in this respect, is the experimental study conducted by Gabriel on 18 speakers of Peninsular Spanish \parencite{gabriel2007fokus} and 50 speakers of Argentinian Spanish \parencite{gabriel2010focus}. Based on elicitation tasks through images, Gabriel came to the following conclusions: speakers  of  Spanish  cross-dialectally prefer  preverbally focused subjects  in  clauses  with  a  full  object NP (\ref{ex27}), while  the postverbal position is favored when the object is realized as a clitic (\ref{ex28}), and in constructions with intransitive verbs, i.e., unaccusatives (\ref{ex29}) or unergatives.

 \ea \label{ex27}
{[\textit{María}]\textsubscript{\textsc{foc}}} \normalsize \textit{comió la manzana}.\\
`María ate the apple.'

\z

 \ea \label{ex28}
\textit{La comió }{[\textit{María}]\textsubscript{\textsc{foc}}}\normalsize.\\
`María ate it.'

\z

 \ea \label{ex29}
 \textit{ Llegaron} {[\textit{Julia y María}]\textsubscript{\textsc{foc}}}\normalsize.\\
`Julia and María arrived.'\\

\z


It becomes clear from these findings that subject focus in Spanish is also regulated by factors such as the argument structure of the verb and the morphological forms of the constituents in the sentence. Given that movement is a costly operation \parencite{chomsky1991linguistics} and hence, when possible, avoided by speakers, in a transitive sentence with two phonetically-expressed arguments, a “light” (e.g., a pronoun) subject and a full object NP, it is easier for speakers to focus the subject \textit{in situ} (by assigning prosodic prominence), than to move it to a postverbal position, such as in (\ref{ex27}). Conversely, when the object is “light”, because it is realized as a clitic, the movement of the subject to a postverbal position is preferred, such as in (\ref{ex28}). The author concludes that ``the mechanisms of syntactic focus marking, i.e., the use of a particular syntactic construction in a given pragmatic context, is governed by strict rules to a lesser extent than suggested in much of the literature” \parencite[189]{gabriel2010focus}.

These findings are in line with Lambrecht's (2010) Preferred Clause Construction (PCC) theory. Based on analyses of corpora of spoken French, the author found an overwhelming preference for a certain sentence type which is defined as ``Preferred Clause Construction” (PCC). In  this construction type,  the  preverbal  position is occupied  by a clitic pronoun and the postverbal position by an XP (typically a single one) with focus value. The PCC is instantiated in 95 to 97\% of all clauses analyzed in the corpora. The results of Gabriel’s study would then prove that this pattern is also valid for Spanish.

As reported by \textcite{garcia2018introduction}, \textcite{zubizarreta2016nuclear} does take the empirical contribution by \textcite{gabriel2010focus} on Argentinean Spanish into account. However, this phenomenon is often dismissed as mere dialectal variation, which does not shed new light on the actual controversy.

Other empirical studies, based on a contrastive elicitation study comprising four European Spanish varieties, namely Canarian Spanish, Castilian Spanish, Basque Spanish (with L1 Basque), and Basque Spanish (with L1 Spanish) have been conducted by \textcite{del2018language}, who, based on 8 participants, observed a number of further patterns of syntactic variation related to focus marking. With regard to information focus on the subject, for example, they found that while speakers of Castilian and Basque Spanish (both with L1 Spanish and L1 Basque) tend to use \textit{in situ} marking, speakers of Canarian Spanish show a significant preference for cleft constructions. Moreover, the  data  shows that, in most  varieties, focused subjects tend to be realized \textit{in situ} regardless of focus type. 

Indeed, clefts constructions have been reported as focus-marking strategies for Spanish as well. The standard assumption in many theoretical studies is that clefts in Spanish are used to mark mainly contrastive focus on the subject (e.g., \cite{lambrecht1996information}). The rationale behind this claim is formulated in terms of economy of the language. Re-ordering of the constituents in the sentence is, by definition, less costly than using a “heavy” syntactic construction. Therefore, if Spanish speakers have more economic means to focus the subject in neutral contexts (i.e., information focus), there is no reason why they would opt for a heavier and costlier strategy. Such claims, however, have been further challenged by empirical evidence.

\textcite{cabrera1999funciones}, for instance, states that simple clefts on the one hand, and (inverted) pseudo-clefts on the other, have different information structural properties, and are not limited to conveying contrastive/corrective interpretations of the subject, as in (\ref{ex30}):

\ea \label{ex30}
\textit{Es} {[\textit{Juan}]\textsubscript{\textsc{foc}}} \normalsize \textit{quien ha llegado tarde}.\\
`It is Juan who arrived late.' \hfill \parencite[4247]{cabrera1999funciones}  
\z


\textcite{feldhausen2015oraciones} have tested these claims through elicitation tasks with 4 speakers of Spanish from Madrid. Their results showed that cleft constructions are used for both focus types, even though certain preferences can be observed: (inverted) pseudo-clefts are favored for neutral focus, while simple clefts are preferred for contrastive focus. After analyzing speakers’ answers prosodically, the authors claim that focus does not always have to bear sentential stress: in clefts, prosodic alignment (phrasing) can be a reliable correlate of focus. In other words, since clefting the subject allows the speaker to position it at the edge of a phonological phrase, prosodic prominence is already achieved by virtue of this position, and assigning pitch accent is no longer necessary.

As for contrastive foci, while theoretical studies claim that fronting (or in the case of the subject, \textit{in situ }focusing) is a formal means of encoding contrastive focus in Spanish (cf. \cite{zubizarreta1998prosody}), other views suggest that contrast can be encoded by other alternative means. \textcite{brunetti2009italian}, for instance, gives  examples  of  contrastive  foci  which  are  not fronted, as well as examples where the contrastive focus is fronted. 
Empirical studies, thus, suggest that variation can be found both in the case of information and contrastive subject focus, and that such variation can be accounted for by considering not just syntactic or exclusively prosodic factors, but also semantic and dialectal ones.

\textcite{garcia2018introduction}, reviewing the question of subject focus in Spanish, point out the frequent lack of comparability of the results issued from the different studies due to i) methodological differences and ii) conceptual or terminological differences regarding the different focus types/categories. The authors stress the fact that many of the claims advanced in the literature are based on methods of data collection that differ considerably (e.g., elicitation tasks, reading tasks, corpora, etc.), and that such differences might have repercussions on the findings.
\textcite{uth2014spanish}, for instance, compares different elicitation tasks, and shows that even slight differences in the design of picture-based elicitation materials may have a crucial impact on both word order and intonation.

Moreover, the classification of focus categories such as ``broad”, ``presentational”, ``neutral”, \say{narrow}, \say{contrastive}, \say{emphatic}, \say{identificational}, \say{exhaustive}, \say{corrective}, \say{mirative}, etc. adds confusion to the picture, since some authors might have a definition of information focus that differs from that of other authors.

Similarly, \textcite{vidal2019nota} present arguments in favor of \textcite{zubizarreta1998prosody}, according to which subject foci that appear in a non-final position cannot be instances of information focus, and have to be considered occurrences of some other focus type (e.g., contrastive/corrective). Their claims are once again based on inconsistencies found at the level of the methodology or the notions used by previous empirical studies. 

Reviewing the method employed by \textcite{gabriel2007fokus}, based on elicitation pictures, the authors agree with \textcite{uth2014spanish}, who claims that the focalized subject in this experiment tend to appear preverbally because it carries the wrong (non-intended) pragmatic inferences: since the answer that participants have to provide is too obvious, it cannot be taken as an instance of new information, but rather as an indication of epistemicity or evidentiality. \textcite{vidal2019nota} therefore advocate caution in the interpretation of these findings. 

In sum, while some authors state that there is no one-to-one correlation between focus type and subject position in Spanish, others reinforce the position taken by \textcite{zubizarreta1998prosody}, and claim that a subject expressing information focus in Spanish must compulsorily appear in the postverbal slot.
 
 \section{Subject focus in French: previous studies}
In the literature on focus in French, it is possible to observe contrastive views as well. Theoretical approaches claim that subject focus can only be expressed by a \textit{c'est} cleft and that, due to the rigid melodic structure of the language, prosodic marking through pitch accent (e.g., \textit{in situ} focusing) is not acceptable (e.g., \cite{lambrecht1996information}). Some empirical studies seem to confirm this claim, whereas others argue that \textit{in situ} strategies are possible through intonational phrasing of the subject constituent (e.g., \cite{fery2001focus}). 

Indeed, the literature on French focus marking conventionally states that French marks focus via syntactic means. Compared to other Romance
 languages or  to  English,  French  is  often  assumed  to  require  special syntactic constructions to mark focus, and seems to be known as the ``black sheep” for not having a flexible prosody.  Yet, over the last twenty years, much work has been done  on  the  prosodic  marking of  information  structure  in  French \parencite{jun2000phonological, fery2001focus, beyssade2015prosodic}.
 Moreover, recent corpus studies have strengthened the argument that a \textit{c'est} cleft is not the only means to mark information focus on the subject, and that \textit{il y a} clefts can also carry out the same IS-function \parencite{karssenberg2017oiseaux, Karssenberg2018}. 
 
Overall, the literature on French reveals a picture similar to the one observed for Spanish: while some authors are in favor of a one-to-one mapping between function and form, others dispute this claim, and show that the same function can be expressed by different, competing means. These studies are discussed in the following sections. Sections 2.3.1 and 2.3.2 summarize research on the realization of French focused subjects from syntactic and prosodic perspectives, respectively. Finally, Section 2.4 illustrates how the present book aims to address the gap in the investigation of this phenomenon for both Spanish and French.

\subsection{Clefts as focalization strategies} 
Cleft sentences are generally assumed to be  focus-marking strategies, although their distribution and IS-properties vary across languages. French, for instance, is a language that is claimed to make a much more abundant use of these constructions, compared to other Romance languages \parencite{katz2000categories, dufter2009clefting, decesare2016sintassi, garassino2022pragmatic_interference}.

Importantly, not all occurrences of clefts are focus-marking devices. Several authors have pointed out that, even in French, a cleft sentence can sometimes be used in different contextual situations which are not related to focus \parencite{prince1978comparison, hedberg1990discourse, delin1992properties, dufter2009focus, mertens2012prosodie}.
Such cases are out of the scope of the present work, and will not be discussed here. This being a study on focus, I will limit the discussion to cleft sentences with a focus-background articulation.

The ``cleft” construction, a term coined by \textcite{jespersen1927modern}, is a bi-clausal construction consisting of a main copular clause, or “clefted element”, and a non-restrictive relative clause that expresses a proposition which could otherwise be conveyed in one single, non-clefted clause. The latter is considered equivalent to the cleft with respect to its truth conditional value and its illocutionary force \parencite[457]{lambrecht2001framework}. The term cleft has since been extended to different types of bi-clausal constructions, such as \textit{c'est} clefts (or it-clefts) pseudo-clefts, or \textit{il y a} clefts (or there-clefts):


\ea \label{ex31} \textbf{\textit{C'est} cleft}\\
\textit{C'est} {[\textit{Marie}]\textsubscript{\textsc{foc}}} \normalsize \textit{qui a acheté le gâteau}.\\
`It is Marie who bought the cake.'

\z

\ea \label{ex32} \textbf{Pseudo cleft}\\
\textit{Ce qui m'embête c'est} {[\textit{qu'ils ne nous donnent pas d'informations}]\textsubscript{\textsc{foc}}}\normalsize.\\
`What annoys me it's that they don't give us any information.'

\z

\ea \label{ex33}\textbf{\textit{Il y a} cleft}\\
\textit{Il y a} {[\textit{une vieille dame}]\textsubscript{\textsc{foc}}} \normalsize \textit{qui habite pas loin d'ici}.\\
`There is an old lady who lives not far from here.'

\z

Clefts have been analyzed from both derivational and functional approaches. Within cartographic syntax, for instance, it has been proposed that all clefted constituents are moved into a designated focus projection in the left periphery \parencite{kiss1998identificational, kiss1999english}. In \textcite{belletti2008answering, belletti2008cp, belletti2012revisiting}, the author proposes a cartographic account in which she argues for the existence of two different focus positions in the clausal map: i) the vP peripheral low focus position in the matrix clause, in the case of new information subject clefts and ii) the high left peripheral position in the reduced CP complement of the copula, in the case of contrastive/corrective object clefts – and also in subject clefts when they are interpreted/used contrastively/correctively \parencite[96]{belletti2012revisiting}. Therefore, under this view, the interpretive effects associated with clefts derive directly from the syntactic position in which the focus is realized (although \cite{haegeman2014architecture} argue against analyzing clefts as occupying a ``high” CP position).

In non-derivational frameworks, on the other hand, clefts are discussed as grammatical constructions in the sense of Construction Grammar \parencite{fillmore1988mechanisms}, i.e., as grammatical  configurations  in  which  morphosyntactic  and prosodic features concur to produce particular form-function-use correspondences.

\textcite{lambrecht1996information} claims that the focus meaning of the two clauses contained in a cleft is non compositional, in the sense that it is not the computable sum of the meaning of its parts. Rather, it is a property of the complex grammatical construction as a whole. Thus, neither its form nor its interpretation can be predicted on the basis of general syntactic and semantic properties of the grammar. Clefts, in this sense, would be grammaticalized constructions whose function is to convey the pragmatic interpretation of focus. 

The author also states that a \textit{c’est} cleft would be an optimal strategy to express subject focus because of a constraint allowing only topical material in preverbal subject position in French. Such constraint would lead to an “increased necessity to resort to the use of clefts” \parencite[490]{lambrecht2001framework}. Since preverbal subjects are generally topical \parencite{chafe1976givenness, keenan1976foregrounding}, an agent receiving focal status is often expressed as an object (i.e., in the focus phrase following the copula clause of the \textit{c’est} cleft). In this position, it does not violate the language prosodic and syntactic constraints that demand that main stress occur at the right-edge of a intonational phrase \parencite{jun2000phonological}. 

Thus, subjects are moved into a dedicated focus position (the clefted element) in which they are interpreted as a focus, and marked as such. The basic premise for this proposal is the observation that languages tend to reserve the preverbal position for topical material – material which is more readily accessible in the discourse. Moreover, this idea is implemented in \textcite{aissen1999markedness} through the harmonic alignment of two scales, (i) a relational scale (Subject $>$ Non-subject) and (ii) a focus scale (Topic $>$ Focus), resulting in the generalization that foci in preverbal subject position are not harmonic and hence should be avoided. 

Therefore, a \textit{c’est} cleft would be an optimal strategy because it allows the desired pragmatic effect of focus on the subject and, at the same time, respects the prosodic requirement of the language (phrase-final stress).
By creating two separate IPs, the cleft allows the  focused  element  to  receive  main  stress by virtue of being situated in a boundary tone, and  to  fulfill  the  rightmost  preference  of the language for accent placement. 

As reported by \textcite{dufter2009clefting}, in the generative tradition \parencite{szabolcsi1981semantics, kiss1998identificational, gussenhoven2008types}, the focus marking function of clefts is considered an epiphenomenon, while all occurrences of clefts are, in fact, strategies to enforce an exhaustive interpretation on the clefted constituent. In other words, the value provided by the  clefted element excludes all the other possible alternatives that satisfy the wh-variable contained in the (implicit or explicit) QUD. This implies that examples like the one found in \textcite[88-89]{jespersen1927modern}, `It is the wife who decides', must be considered semantically equivalent to `Only the wife decides'. 

Similarly, \textcite{lambrecht2001framework, katz1997syntactic} mention that one of the discourse functions of the \textit{c’est} cleft is to identify the X as having the property P, carrying the inference that nothing else in the context displays the property. In other words, a \textit{c’est }cleft would carry semantic exhaustivity encoded in its syntactic form.

A different view is proposed by \textcite{horn1981exhaustiveness}, who suggests that clefts only conversationally imply exhaustivity.
A number of recent experimental studies \parencite{destruel2018interpretation, de2018experimental} has confirmed Horn’s assumption by showing that adding a continuation that cancels the initial exhaustivity of a \textit{c’est} cleft does not affect the truth-conditional value of the sentence. Moreover, it is judged as felicitous by speakers. These results support the prediction that cleft sentences are associated with an exhaustive inference that is cancelable, and therefore not part of their at-issue meaning. \textcite{dufter2009clefting} himself argues that exhaustivity in these constructions should not be taken as an invariant aspect of cleft meaning, because it is not semantically entailed, but pragmatically conveyed and, hence, cancelable. Such argumentation would also justify why occurrences of \textit{c’est} clefts are found in combination with exhaustivity particles, such as \textit{seulement} ‘only’\footnote{As \textcite{dufter2009clefting} rightfully argues, if  exhaustivity were  indeed  part  of  the  semantic  content  expressed  by  the  cleft, one might wonder why restrictive particles such as \textit{only} occur as modifiers of clefted constituents at all. In theory, expressions  with  restrictive semantics cannot be combined felicitously (cf. ?only exclusively the wife).}  or, conversely, with additive particles such as \textit{aussi} ‘too’, which are clear indicators of the non-exhaustive interpretation of the proposition. 

In a similar vein, \textcite{de2015status}, based on a corpus of authentic cleft occurrences collected from different written sources, show that  this meaning component associated with clefts in English and Italian is best accounted for in terms of a conventionalized conversational implicature. Thus, whether exhaustivity is part of the coded meaning, or rather a generalized conversational implicature, is still a matter of debate. Consider, for instance, the following examples taken directly from the corpus under investigation (\textit{sgs} corpus, \cite{adli2011gradient}):

\ea \label{ex34}
 \textit{C'est seulement} {[\textit{moi}]\textsubscript{\textsc{foc}}} \normalsize \textit{ensuite qui suis venue frapper pour la prévenir}.\\
`It is only me, afterwards, who came to warn her.' \hfill (sgs French, 38/187)   
\z


If the exhaustive contribution of a \textit{c'est} cleft is debatable, its IS-function as a focalizing means for the subject is less controversial.
\textcite{destruel2016focus}, for instance, compared spoken colloquial French (CoF) to Standard French (StF) in two elicitation tasks, testing 21 and 21 participants respectively. Her findings show that, in the case of CoF, participants used significantly more clefts when focusing a subject and significantly more canonical sentences when focusing a non-subject. Conversely, in the case of StF, participants  allow grammatical subjects to be focused \textit{in situ} as well as in a cleft \parencite{destruelFery2020_prominence_dualfocus, destruelLalandeChen2024_prosodic_focus}. Moreover, the difference between the two forms was negligible, thus confirming Lambrecht’s (\citeyear{lambrecht1996information}) assumption that the restriction on subject focus occurring in preverbal position is specific to the colloquial grammar of French.
Additionally, she found a slight effect of focus type: participants were found to only produce slightly more clefts when expressing a correction, compared to when marking information focus.

\textcite{reichle2014cleft} combines corpus data (\textit{Europarl} corpus, \cite{koehn2005europarl}) with processing tasks (21 participants), showing that subject clefts are not only more frequent, but also more easily processed than object clefts. He considers frequency and processing costs as indicators of the fact that a \textit{c’est} cleft is the ideal means to express focus on a subject.

Thus, these studies seem to support Lambrecht’s (\citeyear{lambrecht1996information}) claim  that a \textit{c’est} cleft is the preferred way of expressing subject focus in French.
Lambrecht even argues that French is particularly transparent in the mapping between IS onto syntax, and posits a one-to-one correlation between cleft type and IS-function: \textit{c’est} clefts are in fact considered the focus-marking device for argument focus, in particular subject focus. In contrast, \textit{il y a} clefts would express sentence focus, where no information-structural partition is observable, thus conveying an all-new interpretation, with both the clefted element and the relative clause in focus (\ref{ex35}). 

\ea \label{ex35}
 \ {[\textit{Il y a Jean qui a téléphoné}]\textsubscript{\textsc{foc}}}\normalsize.\\
`Jean called.' \hfill \parencite[136]{lambrecht1988presentational}
\z


Moreover, \textcite{lambrecht2010constraints} notes a correlation in terms of information status of the subject constituent: given that \textit{il y a} clefts introduce new referents in the discourse, subjects of these constructions generally appear as indefinite NPs (e.g., \textit{Il y a une vieille dame qui arrive} `There is an old lady who is coming'). \\
Crucially, such a categorical relation between focus type and grammatical construction has been questioned by empirical research. A number of studies has shown that the same structure can correspond to different IS-interpretations, and that the one-to-one correlation between form and function is, if anything, a theoretical assumption more than a proven fact. This also implies that the same IS-function can be expressed by different forms. For instance, it has been shown that \textit{c’est} clefts can also be found in presentative contexts (all-new sentence) or exhibiting a topic-comment partition, where the clefted element is actually the topic and the relative clause is the comment \parencite{dufter2009clefting}.
Similarly, recent corpus studies have shown that \textit{il y a} clefts can occur with a focus-background partition, and not exclusively as all-focus strategies \parencite{skopeteas2010focus, karssenberg2017oiseaux, karssenberg20188}. \\
\textcite{karssenberg2017oiseaux} and \textcite{Karssenberg2018}, for instance, conducting a corpus study on \textit{il y a} clefts, found that a considerable number of these constructions were used to mark focus on the subject, and that the relative clause contained presupposed (old) rather than new information.
To account for this finding, the author argues that what seems to determine the alternation between the two constructions is exhaustivity (or the lack of it), claiming that there is a frequent correlation between \textit{c’est} clefts and exhaustive focus interpretation, while \textit{il y a} clefts are more often found in non-exhaustive contexts. She also adds, however, that exhaustivity is not entailed semantically in the two clefts types, but rather conveyed by contextual cues and is, therefore, cancelable.

Additionally, several prosodic investigations of focus have contributed to weaken Lambrecht’s (\citeyear{lambrecht1996information}) claim  of a deterministic form-function relation in spoken French. These studies are discussed in the next section.

 \subsection{Prosodic realization of subject focus in French}
Despite Lambrecht’s (\citeyear{lambrecht1996information}) claim about focused subjects being obligatorily marked by clefts, some researchers take a canonical sentence like (\ref{ex36}) to be completely acceptable, and argue that subject focus can also be marked via prosodic means \parencite{beyssade2004prosody}.

 \ea \label{ex36}
\ea  \textit{Qui a acheté une mandoline?}\\
 `Who bought a mandolin?'
\ex  {[\textit{Mallarmé}]\textsubscript{\textsc{foc}}} \normalsize \textit{a acheté une mandoline}.\\
`Mallarmé bought a mandolin.' \hfill \parencite[478]{beyssade2004prosody}   
\z \z 


\textcite{belletti2010structures} points out that French speakers do not totally exclude SV answers similar to the English type, even if she does not provide empirical evidence for this claim. Similarly, \textcite{skopeteas2010focus}  found that 10 native speakers produce narrowly-focused subjects in  their canonical sentence-initial position at a high rate (45.3\% of the time). Finally, \textcite{samek2005prosody} develops an optimality-theory account where \textit{in situ} realization is claimed to be the optimal output for subject focus. 

A prominent contribution, in this respect, has been made by \textcite{fery2001focus} who conducted a production task in which 10 native speakers read the answering component of a question-answer pair. Although given the option to change the wording if judged unnatural, participants marked focus on grammatical subjects via prosodic phrasing (i.e., by realizing the subject in its own prosodic phrase), retaining the canonical order in most cases. These findings led her to formulate the Phrasing Hypothesis \parencite{fery2001focus}, according to which prosodic marking of focused subjects in French is possible by creating phrasing on the subject constituent. Therefore, even if assigning pitch accent on a single syllable is banned, the \textit{in situ} strategy is still acceptable by virtue of situating the subject constituent in a boundary tone, thus giving it prosodic prominence. 

Of course, one limitation of this study is that the participants are given scripts for the sentences to produce. One reason why they chose not to vary the sentence structure, and produce a cleft instead, could simply be that they considered changing the word order more costly that maintaining it. Yet, her study showed that prosody is sensitive to information structure in French, a finding that has important implications on the debate about subject focus and its realization.

More recently, \textcite{greif2021correction} conducted an experiment with 16 native speakers, examining whether canonical sentences and cleft constructions could be used with different prosodic realizations with different treatments of corrective focus in typologically different languages, including French. Their findings reveal that both canonical sentences, with a higher prosodic prominence on the edge-tone syllable of the subject, and cleft sentences are possible. When comparing the production task with the results from an acceptability judgement, they found that canonical sentences are slightly more felicitous in the case of object focus than in case of subject focus, but are nevertheless accepted, even when focusing the subject.

This shows that the expression of subject focus in French may be more complicated than previously acknowledged and that more thorough empirical work is still needed to understand the exact conditions under which clefting is required, which cleft types are used as focus-marking strategies and how often the \textit{in situ} strategy occurs. Even if all the strategies discussed above are possible, it is not clear with which distribution they occur and, especially, what determines the optionality among them.

\section{Filling the gap}
 
Considering the studies discussed above, it is safe to conclude that the question of subject focus for both languages is far from being resolved, and that more research is needed. Although many theoretical studies hint at a one one-to-one correlation between form and function, empirical data seems to suggest that such a claim is too strong. 

\subsection{The importance of argument structure}
The overall impression is that Spanish constituent order  is  indeed  free, and  that  focused  constituents  can  appear  in  various  positions, irrespective  of  the type  of  focus  (contrastive or  informational). However, as Gabriel’s (\citeyear{gabriel2007fokus}) study shows, other linguistic factors may interact with focus (and focus types), such as the argument structure of the verb, the morphological form (and hence the information status) of the subject or constituent weight.\footnote{For instance, focus realization might be subject to the ``End-Weight Principle” \parencite{quirk1985comprehensive,hawkins1994performance,hawkins2004efficiency} according to which informationally ``heavy” constituents are placed toward the end of the clause. This principle provides a processing-based explanation for the alignment of focus with clause-final or prosodically prominent positions.} 

A similar scenario can be observed for the use of different cleft formats in French. For instance, it had been initially claimed that sentences like \textit{c’est} clefts or pseudo clefts display the same communicative properties and, as such, are totally interchangeable \parencite{akmajian1970deriving}. It has become clear, however, that the variation among these forms is not totally free and unconstrained, but rather regulated by specific factors. In fact, \textit{c'est} clefts have been claimed to be “preferred with shorter constituents such as pro-forms, whereas pseudo-clefts more readily accommodate longer, in particular clausal, elements in the cleft  position” \parencite[24]{dufter201614}. This implies that the morphological form of the subject constituent (and thus its level of referentiality) has an impact on the realization of focused subjects in French as well. Similarly, to account for the alternation between \textit{c’est} and \textit{il y a} clefts, \textcite{karssenberg2017oiseaux} points out the role of exhaustivity (and hence of focus type). Her assumption, however, was not tested within her study.

If focus type and activation of the subject constituent have been claimed to affect the realization of subject focus in both languages, the role of argument structure has only been investigated in relation to focus in Spanish, i.e., in correlation to Spanish pre- and postverbal subjects. How argument structure interacts with French focus realization, however, remains so far unexplored. In other words, there is no systematic study that investigates whether the argument structure of the verb may affect, say, the choice among different cleft types.

\subsection{Elliptical strategies}
A further limitation may be found in many elicitation tasks, where participants are usually instructed to respond to the stimulus questions with full answers, thus limiting the range of possibilities in the answer type and the spontaneity of expression. These studies tend to exclude any type of elliptical utterance from their analysis (e.g., \cite{skopeteas2010focus}). Since many studies on focus in Romance languages concentrate on syntactic variation (and on the position of the arguments with respect to the verb), the requirement for producing full sentences is not surprising. 

Nevertheless, these elliptical structures are extremely frequent in spoken speech and, at least in dialogical situations, they produce a much more natural effect than their full counterparts. Evidence from studies on register and utterance type suggests that fragmental sentences are not only very frequent, but constitute even the majority of sentence types in spontaneous dialogues. For instance, according to \textcite{du1987discourse}, while speakers tend to produce heavier structures in narratives (e.g., in classical sociolinguistic interviews), full sentences are sometimes unnecessary and redundant in dialogues. In the latter situation, a shorter answer type is perfectly acceptable for the communicational needs of the speakers.

Similarly, \textcite{krifka2008basic} notes that in spontaneous speech, speakers tend to omit the presuppositional part of the information, and to utter only the informative part, i.e., the focus. The same has been claimed by \textcite{fiedler2006subject} in their cross-linguistic investigation of subject focus.
Also according to \textcite{lahousse2012word}, sentences like (\ref{ex37}) may not sound totally natural to certain speakers because the answer to the same questions in spontaneous speech is likely to be the subject alone (as in (\ref{ex38})):

 \ea \label{ex37}
\ea  \textit{Qui a mangé les gâteaux?}\\
 `Who ate the cakes?'
\ex \textit{Ont mangé les gâteaux} {[\textit{Marie, Pierre et Stéphanie}]\textsubscript{\textsc{foc}}}\normalsize.\\
`Marie, Pierre and Stephanie are those who ate the cakes.'\\

\z \z 

\ea \label{ex38}
\ea \textit{Qui a mangé les gâteaux?}\\
 `Who ate the cakes?'
\ex {[\textit{Marie, Pierre et Stéphanie}]\textsubscript{\textsc{foc}}}\normalsize.\\
`Marie, Pierre and Stephanie.' \hfill \parencite[398]{lahousse2012word}   
\z \z 


\textcite{ortega2016focus} points out that the literature on focalization processes in Spanish and beyond rarely considers the interaction between focus and ellipses, even though a considerable number of studies has noted the need for the element surviving ellipses, also known as ``the remnant”,  to be focused and, as such, to escape the ellipsis site (see \cite{merchant1999syntax, merchant1999syntax}). However, the interplay between ellipses and focus only figures marginally in the literature on focalization in Spanish, in spite of the already-noted connection between these structures and their IS-properties (\cite{merchant1999syntax}, \textit{inter alia}). 

Finally, \textcite{vidal2019nota}, reviewing the studies on the realization of subject focus in Spanish, stress that in wh-questions eliciting information focus on the subject, the preferred answer is always a fragment and not a complete one (also noted by \cite{fery2017intonation}). The authors claim that the reason lies in the fact that the default option for presupposed (backgrounded) information is to be elided. Thus, when participants are forced to produce sentences including all the elements in the utterance, ``they are prevented from producing a natural, spontaneous answer” \parencite[9]{vidal2019nota}.

As for the investigation of focus in French, several authors have noted that the most natural answer to a subject interrogative in dialogues is what has been defined in the literature as “reduced clefts” \parencite{belletti2005answering, belletti2008cp, belletti2008answering}.\footnote{\textcite[538]{doetjes2004cleft} refer to such instances as ``truncated clefts” (see also \cite{buring1998identity, hedberg2000referential}).} 

 \ea \label{ex39}
\ea  \textit{Qui} (\textit{est-ce que qui}) \textit{a parlé?}\\
 `Who spoke?'
\ex \textit{C'est} {[\textit{Jean}]\textsubscript{\textsc{foc}}} \normalsize (\textit{qui a parlé}).\\
`It is Jean (who spoke).' \hfill \parencite[192]{belletti2008cp}    
\z \z


This term indicates the realization of a cleft without the presupposed part of the sentence \parencite{frascarelli2013pseudo}, and has often been the object of discussion for French and Brasilian Portuguese \parencite{lobo20172}, even though it can be found in other languages, like Spanish.

\ea \label{ex40}
\ea  \textit{Quién ha hablado?}\\
 `who spoke?'
\ex \textit{Ha sido} {[\textit{Juan}]\textsubscript{\textsc{foc}}} \normalsize (\textit{quien ha hablado}).\\
`It was Juan (who spoke).'

\z \z


\textcite{karssenberg2017oiseaux} notes that, in the case of French \textit{il y a} clefts as well, the relative clause often tends to be omitted by speakers when it contains discourse-given information (presuppositional information), leaving only the clefted element introduced by \textit{il y a}, such as in (37b):


\ea \label{ex41}
\ea  \textit{Il n’y a personne qui monte}.\\
 `There's nobody going upstairs?'
\ex \textit{Si, il y a} {[\textit{Pierre}]\textsubscript{\textsc{foc}}} \normalsize (\textit{qui monte l’escalier}).\\
`There is Pierre (who goes upstairs).' \hfill \parencite[202]{karssenberg2017oiseaux}    
\z \z


Evidence from acquisition \parencite{lobo2019effects} shows that the production of full clefts increases significantly with age, whereas the proportion of reduced clefts, which are more frequent than full clefts, decreases from the age of 4. \textcite{lahousse2019emergence} conducted a corpus study on the spontaneous speech of early childhood (age 1-3 y.o.), and found that reduced clefts are more frequent than other cleft types, and that full clefts tend to appear at a later stage of this age range. 

As for the usage of these structures, \textcite{belletti2005answering} claims that reduced clefts are preferred as answers in contexts where the subject interrogative immediately precedes the answer as in (\ref{ex38}), even if her claim is not based on empirical evidence. Similarly, \textcite{reichle2014cleft} claims that in Q-A pairs with fast turns, ``full” clefts might be considered unnecessarily wordy, and the presupposition contained in the relative clause is not necessarily repeated in the response. He states: “Such abbreviated responses are possible when the clefted relative clause is highly activated due to having just been mentioned” \parencite[8]{reichle2014cleft}. However, the longer responses (with a full repetition of the cleft in the response) are possible, and findings from judgment tasks indicate that native speakers find such pairings acceptable \parencite{reichle2014cleft}.

\textcite{katz2000categories} had already conducted an acceptability judgement study, comparing the participants' evaluation of full clefts vs. reduced clefts. The author concludes that both structures are judged to be felicitous by speakers of French. Apart from this acceptability task, both in the case of reduced clefts and in the case of fragments  there are no systematic empirical studies that investigates the frequency of these strategies (either reduced clefts or fragments) in the two languages, nor the pragmatic contexts in which they occur.

The need for more empirical investigation is also advocated by \textcite{destruel2016focus}, who points out that: 

“\textit{Full sentences are not typically used in natural discourse, and ensuring that participants use such forms may have led to elicited data that contains a biased, non-optimal variant of participants' speech. This is indeed potential drawback that will need to be addressed in future work. The present experiment constitutes a first step in the exploration of focus-marking strategies in different varieties of French that will certainly benefit from further empirical work exploring more naturally-occurring data (i.e., corpus studies and other elicitation methods like semi-directed interviews) to corroborate the current findings}.” \parencite[309]{destruel2016focus}.

It is important to point out that there is an ongoing debate on the exact syntactic nature of structures such as fragmental answers or reduced clefts. 
Concerning fragments, some  authors  consider these structures as covertly clausal and elliptical \parencite{merchant2005fragments}, while others claim that they are in fact base-generated constituents \parencite{stainton1998quantifier, jonathan2000interrogative, jacobson2016short}, and not the result of some kind of deletion.

Similarly, in the debate concerning reduced clefts, some authors are in favor of the “reduction hypothesis” \parencite{jenkins2017cleft}, which posits a base-derived account, according to which these structures would be parallel in “deep and surface” structure to clefts without the relative clause (i.e., the ``shortened/truncated cleft analysis”). Others, like \textcite{milsark2014existential, birner2007functional} present several arguments against such a truncation (or reduction) hypothesis, concluding that the similarities between ‘full’ clefts and copular sentences without relative clauses are due to functional overlap, and that one should not be derived from the other \parencite[319]{birner2007functional}. While the first view calls these structures ``reduced” or ``truncated clefts”, the second view refers to them as copular or, in the case of the ones introduced by \textit{il y a}, existential sentences.

\subsection{The contribution of this book}
Determining the exact syntactic nature of these constructions is out of the scope of this book, as the present corpus study is mostly concerned with what these structure do from a functional perspective, and how they contribute to the universe of discourse. Nevertheless, for the purpose of the present investigation, I maintain the term ``fragment” and ``reduced cleft”, and I consider that, from a pragmatic perspective, both fragmental answers and reduced clefts are answering strategies where only the focused element is visible\footnote{I follow \textcite{yoshida2013predictive} and assume that, at least from a communicative perspective, focused material cannot undergo elision.
}, while the presupposition is implicitly shared by speaker and hearer in the CG. In other words, I do not take a position on whether, syntactically, a reduced cleft is originally a full cleft whose relative clause is elided. Rather, I employ the term reduced cleft to distinguish it from a full cleft because I consider that, communicatively, the presuppositional part of information has been reduced or elided by the speaker.

Irrespective of their syntactic nature, a corpus study that investigates the frequency of these strategies compared to other focus-marking strategies is needed, and that is one of the main objective of this book.
Although the studies so far conducted have made an extremely valuable contribution  to the investigation of subject focus, it has to be recognized that a systematic corpus investigation that takes into account both full and elliptical answers would shed light on the phenomenon both from a language-internal and from a comparative perspective. 

In light of what has been discussed above, the main goals of the present work are: i) to identify, through a corpus analysis, which answering strategies (both full and elliptical) are used in French and Spanish to express focused subjects, and to determine their distribution in spontaneous speech, (ii) to investigate the contexts and motivations that determine a speaker's decision to opt either for a full or an elliptical answer in both languages and, finally, (iii) to determine whether the factors such as focus type, activation of the subject constituent and argument structure of the verb have an impact on focus realization in both languages.
